\documentclass[12pt,a4paper]{article}

\input{/home/tabaka/Documents/GitHub/Defs/defs.tex}


\setcounter{zadaniaRozdzial}{8}

\begin{document}
	
	\begin{center}
		\LARGE Zadania na lekcję 8
	\end{center}
	\vspace{1.5cm}
	
	\ZbiorZadanie{Udowodnij twierdzenie sinusów.}
	
	\ZbiorZadanie{Udowodnij twierdzenie cosinusów.}
	
	\ZbiorZadanie{W trójkącie $ABC$ mamy dane: $|AC|=3\sqrt{3}$, $|BC|=3$ oraz $\angle A=30^\circ$. Oblicz miary pozostałych kątów tego trójkąta.}
	
	\ZbiorZadanie{Boki trójkąta $ABC$ mają długości $|AB|=8\,\mathrm{cm}$, $|BC|=6\,\mathrm{cm}$, $|AC|=5\,\mathrm{cm}$. Sprawdź, czy trójkąt jest ostrokątny, prostokątny czy rozwartokątny.}
	
	\ZbiorZadanie{Boki trójkąta $ABC$ mają długości $|AB|=10$, $|BC|=3$, $|AC|=5$. Sprawdź, czy trójkąt jest ostrokątny, prostokątny czy rozwartokątny.}
	
	\ZbiorZadanie{Dwa boki trójkąta $ABC$ mają długości: $|AB|=7\,\mathrm{cm}$, $|BC|=8\,\mathrm{cm}$. Wiedząc, że $\angle ACB=60^\circ$, oblicz długość boku $AC$.}
	
	\ZbiorZadanie{W trójkącie prostokątnym równoramiennym $ABC$, $|AC|=|AB|$, poprowadzono środkowe $CD$ i $BE$, które przecięły się w punkcie $M$. Wykaż, że $\cos\angle DMB=\frac{4}{5}$.}
	
	\begin{center}
		\LARGE Zadania domowe -- lekcja 8
	\end{center}
	\vspace{1.5cm}
	
	\setcounterref{zadaniaIndex}{0}
	
	\ZbiorZadanieTwoXTwo{Oblicz promień okręgu opisanego na trójkącie $ABC$, jeśli:}
	{$|AC|=3$ oraz $\angle CBA=60^\circ$}
	{$|BC|=4$ oraz $\angle BAC=150^\circ$}
	{$\angle A=72^\circ$, $\angle B=63^\circ$, $|AB|=\sqrt{8}$}
	{$\angle B+\angle C=150^\circ$, $|BC|=13$.}
	
	\ZbiorZadanie{W trójkącie $ABC$ mamy dane: $|AB|=12$, $|BC|=6\sqrt{2}$ oraz $\angle A=30^\circ$. Oblicz miary pozostałych kątów tego trójkąta.}
	
	\ZbiorZadanie{Wykaż, że jeśli promień okręgu opisanego na trójkącie jest równy długości jednego z boków trójkąta, to trójkąt leżący naprzeciw tego boku ma miarę $30^\circ$ lub $150^\circ$.}
	
	\ZbiorZadanieTwoXTwo{Rozwiąż trójkąt $ABC$, jeśli dane są:}
	{$|AB|=3$, $|BC|=9$, $\angle ABC=120^\circ$}
	{$|AB|=2\sqrt{2}$, $|BC|=6$, $\angle ABC=45^\circ$}
	{$|AB|=10$, $|BC|=9$, $|AC|=5$}
	{$|AB|=4$, $|BC|=10$, $|AC|=12$.}
	
	\ZbiorZadanie{Dwa boki trójkąta $ABC$ mają długości: $|AC|=10$, $|BC|=6$. Wiedząc, że $\angle ACB=60^\circ$, oblicz długość środkowej $CD$.}
	
	
\end{document}